\chapter*{Abstract}
\addcontentsline{toc}{chapter}{Abstract}

\vspace{0.6cm}

Solar flares are sudden releases of magnetic energy in the Sun's atmosphere
whose radiation and ejecta drive space-weather hazards at Earth. Forecasting
them from the observable state of a source active region remains an open
problem, made hard by the rarity of large events and by the fact that the
governing coronal currents are seen only through their photospheric footprint.
Magnetic \emph{winding flux}---a signed, topological measure of how field lines
wrap around one another, concentrated along polarity inversion lines---has been
proposed as a physically motivated precursor, but existing work reduces the
winding map to a single spatial integral before forecasting, cancelling the
signed, localised structure that carries the signal.

This thesis takes a first step toward using the winding-flux map without that
reduction. A Convolutional LSTM (ConvLSTM) encoder--forecaster is trained to
predict the short-term evolution of the two-dimensional winding field of an
active region from a ten-frame history, under a strict
leave-one-active-region-out protocol. A key empirical finding organises the
results: the stored winding field behaves as a near-white \emph{rate}, whereas
its time-integral---the accumulated total winding---is smooth and strongly
persistent, and it is the total that is the appropriate forecasting target. On
the total field the model reproduces the large-scale evolution of held-out
active regions; its skill is reported honestly against a persistence baseline,
and the roles of the spatial low-pass and of the rate--total distinction are
analysed rather than hidden. The work establishes both a working spatiotemporal
pipeline for winding-flux maps and a clear-eyed account of what such
forecasting can and cannot deliver, motivating the structure-aware direction for
subsequent flare-prediction research.
